\section{The remaining global question, stated as an exact estimate}
\label{sec:targets}

\subsection{A complete analytic implication}
For clarity, we prove the one-sided analytic implication used in the previous section. This is a form of Landau's positivity principle, the same mechanism isolated in the Lean source.

\begin{lemma}[One-sided Mellin exclusion]
\label{lem:mellin}
Let \(T:[1,\infty)\to\mathbf R\) be locally integrable and of polynomial growth. Suppose its Mellin transform
\[
 M(s)=\int_1^\infty T(x)x^{-s-1}\,dx
\]
continues meromorphically to \(\Rea s>\theta\), with only simple poles, and every real pole in that half-plane has strictly positive residue. If
\(T(x)\le Cx^\theta\) eventually, then \(M\) has no pole in \(\Rea s>\theta\).
\end{lemma}

\begin{proof}
For a sufficiently large \(X\), the function \(f(x)=Cx^\theta-T(x)\) is nonnegative on \([X,\infty)\). Its Mellin transform is initially
\[
 Q(s)=\frac{CX^{\theta-s}}{s-\theta}
       -M(s)+\int_1^X T(x)x^{-s-1}\,dx.
\]
The finite integral is entire. Let \(\sigma_c\) be the abscissa of convergence of the nonnegative integral defining \(Q\). Polynomial growth makes it finite above; it may be \(-\infty\).

The Landau fact used here is that a nonnegative Mellin integral with finite abscissa has a singularity at the real point \(\sigma_c\). Briefly, if the transform were holomorphic there, its Taylor expansion from a real point to the right, together with nonnegative moments and monotone convergence, would force convergence slightly to the left of the abscissa, a contradiction.

If \(\sigma_c>\theta\), the displayed meromorphic expression makes that singularity a simple pole with residue equal to the negative of the corresponding residue of \(M\). It is therefore negative. But \(Q(\sigma)\ge0\) for real \(\sigma>\sigma_c\); a negative simple-pole residue would make \(Q(\sigma)\to-\infty\) as \(\sigma\downarrow\sigma_c\), which is impossible. If \(M\) has no real pole there, there is no singularity at all, giving the other contradiction. Hence \(\sigma_c\le\theta\). Absolute convergence then makes \(Q\), and thus \(M\), holomorphic in \(\Rea s>\theta\).
\end{proof}

\begin{theorem}[A concrete nonabelian RH target]
\label{thm:a5target}
Let \(E/\mathbf Q\) be Galois with group \(A_5\), and let \(F=\chi^4\) and \(T_F\) be as in Section~\ref{sec:artin}. Suppose that, for every \(\varepsilon>0\),
\[
 T_F(x)\le C_{E,\varepsilon}x^{1/2+\varepsilon}
 \qquad\text{for all sufficiently large }x.
 \tag{A5}
\]
Then every nontrivial zero of \(\zeta_E\) has real part \(1/2\). Moreover, every individual Artin factor has neither a zero nor a pole off that line inside the critical strip.
\end{theorem}

\begin{proof}
The prime-power coefficient bound gives polynomial growth and local integrability of \(T_F\). Its Mellin transform was computed exactly in the preceding section. The point \(s=1\) is removable. In \(\Rea s>0\), its other poles are simple and occur at zeros of \(L(s,F)\); their residues are
\[
 \frac{\ord_\alpha L(s,F)}{\alpha(1-\alpha)}.
\]
For real \(0<\alpha<1\) these are positive. There are no zeros for real part greater than or equal to \(1\), apart from the already removed pole at \(1\); this follows from the Artin--Hecke product identity and the classical zero-free boundary for the factors.

For each \(1/2<\theta<1\), condition (A5) and Lemma~\ref{lem:mellin} exclude zeros of \(L(s,F)\) with real part greater than \(\theta\). The multiplicity comparison \(d/4\le e_F\) excludes zeros of \(\zeta_E\) there. Letting \(\theta\downarrow1/2\), and using the functional equation of \(\zeta_E\), proves its RH.

At an off-line point in the strip, \(d=\ord\zeta_E=0\). The cyclic-restriction inequality \(|\Theta(g)|\le d\) forces the entire Heilbronn character to be zero, so all the individual Artin orders vanish.
\end{proof}

The last assertion permits possible individual Artin poles on the critical line. It does not prove the Artin holomorphy conjecture. The hypothesis (A5) is the exact missing arithmetic statement; this report does not establish it.

\subsection{Why this is a method transfer, and where it stops}
The fourth-tensor construction supplies a nonnegative prime statistic, meromorphic continuation, unconditional holomorphy of the aggregate, and a residue dictionary without cancellation. Those are enough for the analytic endgame. They do not build the arithmetic object whose determinant comparison proves (A5).

A successful adaptation of the source would still have to construct integral sections or forms for this tensor statistic; show that the short sections span almost all the required space; prove a matching lower covolume bound; control ramification and congruence denominators; and keep every loss within the square-root scale. Tensor multiplication of characters by itself proves none of those estimates.

The general faithful-tensor theorem is useful when a geometric construction naturally produces tensor powers of a representation. It supplies a finite, explicit amplification threshold. The sharp \(A_5\) certificate reduces that threshold from a sixth-power pointwise test to the fourth power. It is not the best choice if the only optimization objective is detector degree: the prior degree-16 construction remains smaller.

\subsection{Assessment of candidate problems}
\begin{center}
\small
\begin{tabularx}{\textwidth}{@{}p{.25\textwidth}XX@{}}
\toprule
Question & Progress actually supplied & Missing input or obstruction\\
\midrule
Optimal exact torsion specialization & Solved for the fixed aggregate objective, including sharpness and bounded series & Separate exact individual constraints are a different and obstructed objective.\\
Lubin--Tate specialization & Exact same minimax theorem with residue cardinality \(q\) & No global conductor or density conclusion follows from local order alone.\\
Nonabelian Dedekind RH & Explicit cancellation-safe analytic criterion and tensor certificate & A new arithmetic proof of (A5).\\
Artin holomorphy & Certain aggregate tensor products are holomorphic without the conjecture & Positivity of aggregate orders does not force each constituent order to be nonnegative on the critical line.\\
Simultaneous finite-order nonvanishing & A sharp local specialization budget for nonzero analytic packets & Integral interpolation, nonzero witnesses, admissible character realization, and conductor counting.\\
BSD or Selmer vanishing & Better aggregate determinant specialization may reduce a local search & Source hypotheses require finite conjugate-dual terms; the existing duality obstruction survives.\\
Vandiver or cyclotomic vanishing & The source supplies finite-level length bounds & A finite positive bound is not zero, and the required comparison with the target group must be proved.\\
\bottomrule
\end{tabularx}
\end{center}

For the Selmer direction, the source's \path{SelmerData.lean} explicitly assumes a finite global congruence term, finite local terms, residual spanning conditions, and an inertial scalar condition. Its general twisting theorem gives a bound of the form \(\beta_T+\varepsilon\). The rational conjugate-duality obstruction established in the preceding report is unaffected by anticyclotomic twisting \cite{structuralprior}. An improved torsion ceiling cannot make an infinite \(\beta_T\) finite.

For the cyclotomic direction, \path{Lib/Cyclotomic/OddSha.lean} bounds particular finite-level cocycle modules by valuations of an explicitly bounded nonzero integer. Even after making the comparison with a conventional cohomology group, one still needs an argument that the relevant valuation vanishes. No such argument is supplied here.

\subsection{Recommended mathematical continuation}
The best concrete next step is to formalize the sharp shell theorem and use it at source call sites whose loss is genuinely aggregate. Its hypotheses are local and its proof is finite enough to be separated from the rest of the RH development. The included formalization plan lists the required lemmas without pretending they have already been checked in Lean.

For a global extension, the decisive experiment is a construction of the upper and lower determinant estimates for one nonabelian test family, with every rank and denominator loss written out. Another character identity alone would add less than resolving even one of those arithmetic estimates. The \(A_5\) tensor formula gives a fully specified statistic against which such a construction can be tested.

